% Appendices A and B (C is zh/app_index.tex, D is zh/app_timeline.tex).
\chapter{数学工具}\label{app:math}

本附录收集全书用到的数学工具，每一条注明在哪一章首次使用。计划包括：
\begin{itemize}
\item 差分方程与稳定性（第 25、44、50 章）；
\item 对数、增长率与常用近似（第 2、3 章）；
\item 拉格朗日乘子与跨期最优化，欧拉方程（第 39、53 章）；
\item 对数线性化（第 53、54 章）；
\item 理性预期模型的求解与 Blanchard–Kahn 条件（第 48、55 章）；
\item 向量自回归入门（第 80 章）。
\end{itemize}

\chapterstub

\chapter{符号与数据来源}\label{app:symbols}

\section{符号约定}

除非另有说明，全书采用下列符号。各章第一次使用时会再说明一次。

\begin{table}[htbp]\centering\small
\caption{全书的主要符号}\label{tab:app-symbols}
\begin{tabularx}{\linewidth}{@{}l Y@{}}
\toprule
符号 & 含义 \\\midrule
$Y$，$y$ & 产出；小写 $y$ 在第 44 章以后多指产出缺口（对数偏离） \\
$P$，$\pi$ & 价格水平；通胀率（每期） \\
$i$，$r$ & 名义利率；实际利率 \\
$r^{*}$ & 中性（自然）实际利率 \\
$g$ & 名义或实际 GDP 增长率（在债务动态中与 $r$ 的口径一致） \\
$b$，$s$ & 政府债务与 GDP 之比；基本盈余与 GDP 之比 \\
$M$，$V$ & 货币存量；货币流通速度 \\
$e$ & 名义汇率，定义为一单位外币折合多少本币（$e$ 上升表示本币贬值） \\
$\E_t$ & 以第 $t$ 期信息为条件的期望 \\
\bottomrule
\end{tabularx}
\end{table}

\section{数据来源}

全书使用的数据来源与 \texttt{data/SOURCES.md} 一致；每一份数据快照的序列代码、下载日期与许可条款记录在快照旁的元数据文件中。

\chapterstub
